\documentclass[11pt]{article}
\usepackage{series}
\usepackage{amsmath,amssymb,amsthm}
\usepackage{geometry}
\usepackage{hyperref}

\geometry{margin=1in}

\title{\textbf{Dynamics of Coherence Capacity:\\ Transport, Concentration, and Exhaustion}}
\author{Peter Nero}
\date{January 2026}

\begin{document}
\maketitle

\begin{abstract}
We develop the dynamical theory of coherence capacity introduced in
``Coherence Capacity as the Fundamental Resource of Effective Physics.''
Coherence capacity is defined as the finite stability margin that allows
a projection-based effective description to remain predictive.
Under minimal axioms, we show that coherence capacity admits a local
transport law and behaves as a conserved resource within admissible
regimes. We prove that capacity transport generically exhibits focusing,
leading to concentration, bottleneck formation, and exhaustion on
codimension-one surfaces. These results provide the dynamical substrate
for gravitational attraction, horizon formation, thermalization, and
irreversibility, without modifying fundamental dynamics or introducing
new degrees of freedom. All transport statements are effective/encoding-level, slab-local, and conditional on regularity and constitutive closure assumptions stated explicitly below.
\end{abstract}

\section{Introduction}

In the preceding work we identified coherence capacity as the structural
resource governing the validity of effective physical descriptions.
Exhaustion of coherence capacity forces noninvertibility of the effective
evolution, producing irreversibility even when the underlying dynamics
is deterministic and invertible.

The present paper extends that analysis from kinematics to dynamics.
Specifically, we address the question: how does coherence capacity move,
redistribute, concentrate, and fail?

We show that coherence capacity is not merely a diagnostic but a physical
resource admitting local transport and conservation laws. Within
admissible regimes, capacity is conserved and redistributed; under
persistent strain it concentrates and forms bottlenecks where effective
descriptions necessarily break down.

This transport picture provides the missing dynamical link between
capacity exhaustion and familiar physical phenomena such as gravitational
attraction, horizon formation, and loss of memory.

\section{Levels of Description}

We distinguish three levels.

\subsection{Upstairs configuration space}

Let $(X,\mathcal B)$ be a measurable configuration space equipped with an
invertible evolution $\Phi:X\to X$ (or a flow $\Phi_t$). No fundamental
time parameter is assumed; $\Phi$ defines only an ordering of
configurations.

\subsection{Projection and observables}

Observable physics is defined by a measurable projection
\[
P:X\to Y,
\]
where $Y$ is the space of effective states. The projection is generally
noninjective.

\subsection{Effective shadow}

The induced observable evolution is
\[
T := P\circ \Phi : X \to Y.
\]
All effective laws refer to this shadow level.

\section{Axioms}

The results of this paper rest on the following minimal axioms.

\begin{description}
\item[Axiom A1 (Invertible fundamental dynamics).]
The evolution $\Phi$ is invertible on $X$.

\item[Axiom A2 (Projection).]
The observable map $P:X\to Y$ is measurable and noninjective.

\item[Axiom A3 (Admissibility).]
There exists a subset $A\subset X$ on which the projection-based
description is stable and predictive.

\item[Axiom A4 (Finite coherence capacity).]
Stability of the effective description on $A$ is controlled by a finite
margin that can be exhausted under disturbance.
\end{description}

No assumption of fundamental stochasticity, collapse, or time asymmetry
is made.

\section{Coherence Capacity}

\subsection{Definition}

\begin{definition}[Coherence capacity]
For $x\in A$, the coherence capacity $\mathcal C(x)$ is the maximal
disturbance amplitude such that the projection-based effective
description remains stable in a neighborhood of $x$.
\end{definition}

$\mathcal C(x)>0$ characterizes admissible regions; $\mathcal C(x)=0$
marks admissibility barriers.

\subsection{Basic properties}

\begin{lemma}
$\mathcal C(x)>0$ for all $x\in A$.
\end{lemma}

\begin{lemma}
If $x_n\in A$ converges to $x\notin A$, then $\mathcal C(x_n)\to 0$.
\end{lemma}

These properties follow directly from the definition of admissibility.

\section{Capacity as a Physical Resource}

We now elevate coherence capacity from a diagnostic to a physical
resource.

\begin{definition}[Coherence capacity resource]
Coherence capacity is a locally defined scalar resource whose preservation
is necessary and sufficient for the continued validity of a projection-
based effective description.
\end{definition}

This resource can be redistributed, concentrated, and exhausted, but not
created from nothing within admissible regimes.
\paragraph{Interpretive clarification.}
Referring to coherence capacity as a ``resource'' is a descriptive choice, not an ontological
claim. Coherence capacity is not a new degree of freedom, field, or conserved charge. The
resource language encodes the fact that admissibility margins can be redistributed, depleted,
or concentrated within an effective description, even though the underlying dynamics on $X$
remains invertible. All statements in this section concern the structure of effective
description, not additional microscopic physics.

\section{Capacity Transport and Currents}

\subsection{Existence of a capacity current}

\begin{theorem}[Existence of a coherence-capacity current]
Within any admissible basin, there exists a locally defined effective current 
$J^\mu_{\mathcal C}$ on $Y$, defined up to
coarse-graining, such that
\[
\nabla_\mu J^\mu_{\mathcal C} = 0
\]
holds pointwise inside the basin. Concretely, take any admissibility margin functional $C$ defined on the coherent basin and
consider its pushforward density under the projection map $P$ restricted to that basin; under
the stated absolute continuity and regularity assumptions, this pushforward admits a local
representation by an effective current.

\end{theorem}

\begin{proof}
Inside an admissible basin the projection remains stable and the
effective description is predictive. The invertibility of $\Phi$ ensures
that admissibility-defining structure is preserved under evolution.
Pushing forward this structure defines a local flow of coherence capacity
on $Y$, yielding a divergence-free current $J^\mu_{\mathcal C}$.
Here ``current'' refers to an effective representation of admissibility transport under the
assumption that Y admits a smooth local chart and that the pushforward of admissibility-
preserving structure is absolutely continuous.

\end{proof}

\subsection{Continuity equation}

The local conservation law
\[
\boxed{
\nabla_\mu J^\mu_{\mathcal C} = 0 \qquad (\mathcal C>0)
}
\]
expresses conservation of coherence capacity within admissible regimes.

\section{Capacity Nonconservation at Admissibility Barriers}

Conservation of coherence capacity holds only within admissible regions.
At admissibility barriers, the defining inequalities of stability fail and
capacity is no longer preserved.

\begin{theorem}[Capacity exhaustion at admissibility barriers]
Let $B\subset Y$ be the projection of an admissibility barrier. Then the
capacity current satisfies
\[
\nabla_\mu J^\mu_{\mathcal C} = \mathcal S_B,
\]
where $\mathcal S_B$ is a distribution supported on $B$.
\end{theorem}

\begin{proof}
At an admissibility barrier $\mathcal C\to 0$ and the projection $P$ loses
regularity. The pushforward flow defining $J^\mu_{\mathcal C}$ ceases to be
well-defined across $B$. As a result, the local conservation law acquires
a source term supported on the barrier.
\end{proof}

\begin{corollary}
All irreversible changes in the effective description correspond to flux
of coherence capacity through admissibility barriers.
\end{corollary}

This establishes irreversibility as a structural necessity rather than a
dynamical assumption.

\section{Capacity Flux Direction and Strain}

The continuity equation alone does not determine the direction of capacity
transport. Physical interpretation requires identifying what drives
capacity redistribution.

\subsection{Capacity strain}

\begin{definition}[Capacity strain]
Let $\mathcal L(x)$ denote the local load imposed on the effective
description by interactions, curvature, entanglement growth, measurement
coupling, or environmental noise. Capacity strain is the tendency of
$\mathcal C$ to decrease in regions where $\mathcal L$ is large.
\end{definition}

\begin{assumption}[Downhill transport principle]
Coherence capacity is transported so as to oppose capacity strain: capacity
flows away from regions of lower effective stability toward regions of
higher effective stability. 

\end{assumption}

This assumption is a constitutive choice describing how admissibility margins are encoded at
the effective level; it does not introduce new microscopic dynamics.
\section{Capacity Focusing and Concentration}

We now show that capacity transport generically leads to concentration.

\begin{definition}[Flow direction and expansion]
Where $J_C^\mu \neq 0$, define the unit flow
direction $u^\mu := J_C^\mu / \|J_C\|$ (with $\|\cdot\|$ taken with respect to the effective
metric on $Y$ when such a chart exists). Define the expansion of the capacity flow by
\[
\theta := \nabla_\mu u^\mu .
\]
Inside admissible basins, $\nabla_\mu J_C^\mu = 0$ expresses conservation of the effective
capacity current, but $\theta$ need not vanish; instead, it encodes the focusing or
defocusing of the integral curves of the normalized flow $u^\mu$.

\end{definition}

\begin{theorem}[Capacity focusing]
Assume persistent capacity strain along integral curves of $u^\mu$. Then
there exists $\kappa>0$ such that
\[
\frac{d\theta}{d\tau} \le -\kappa\,\theta^2
\]
along the flow.
\end{theorem}

\begin{proof}
Capacity strain induces gradients in $\mathcal C$ that bias the flux
according to the downhill transport principle. Contractivity of projection
ensures that deviations from uniform flux amplify concentration rather
than dispersion, producing a quadratic suppression term analogous to
Raychaudhuri focusing.
\end{proof}

\begin{corollary}
If capacity strain persists for sufficient duration, coherence capacity
concentrates and $\theta$ becomes negative in finite parameter distance.
\end{corollary}

\section{Bottleneck and Boundary Formation}

Capacity concentration has a sharp endpoint.

\begin{theorem}[Bottleneck formation]
If inward coherence-capacity flux into a compact region exceeds outward
transport for sufficient duration, then $\mathcal C\to 0$ on a
codimension-one hypersurface surrounding the region.
\end{theorem}

\begin{proof}
By the focusing theorem, persistent negative expansion drives capacity to
zero in finite parameter distance. Conservation forbids disappearance of
capacity elsewhere, so depletion localizes on a bounding hypersurface in the effective transport representation.
\end{proof}

\begin{definition}[Capacity bottleneck]
A capacity bottleneck is a connected hypersurface on which $\mathcal C=0$
and across which the effective description is noninvertible.
\end{definition}

Capacity bottlenecks are the structural origin of horizons and collapse
boundaries.

\section{Sources of Capacity Strain}

The following processes contribute to capacity strain:

\begin{itemize}
\item \textbf{Interaction strength:} strong coupling increases projection
load.
\item \textbf{Curvature:} geometric distortion raises the cost of
maintaining stability.
\item \textbf{Entanglement growth:} rapid correlation buildup strains
coherent truncation.
\item \textbf{Measurement coupling:} external apparatus imposes additional
projection constraints.
\item \textbf{Noise and environment:} uncontrolled coupling drains
capacity through noncoherent channels.
\end{itemize}

Any combination of these can drive $\mathcal C\to 0$.

\section{A Minimal Illustrative Model}

The following model is purely illustrative and does not represent a proposed fundamental
dynamical law for coherence capacity; it serves only to visualize focusing, relaxation, and
barrier formation under admissibility-preserving transport.


Consider a one-dimensional effective space with coordinate $x$, coherence
capacity $\mathcal C(x,t)$, and current
\[
J_{\mathcal C} = -D\,\partial_x \mathcal C + v(x)\,\mathcal C,
\]
where $v(x)$ represents capacity strain and $D>0$ is a contractivity
constant.

The continuity equation becomes
\[
\partial_t \mathcal C
= D\,\partial_x^2 \mathcal C - \partial_x(v\,\mathcal C).
\]

If $v(x)$ is inward-directed and sufficiently strong, solutions develop a
finite-time zero of $\mathcal C$, signaling breakdown of the effective
description at a moving boundary.

This toy model captures relaxation when capacity is abundant, focusing
under strain, and barrier formation without singular dynamics.

\section{Discussion}

Coherence capacity behaves as a genuine physical resource: it is
conserved within admissible regimes, transported and concentrated by
interaction and geometry, and exhausted at moving boundaries. These
properties are independent of microscopic details and follow solely from
projection-based effective description with finite stability margins.

In subsequent work we will show that capacity concentration produces
gravitational attraction, capacity bottlenecks generate horizons and area
laws, and competing transport regimes yield thermalization or persistent
memory.

\section{Conclusion}

We have developed a dynamical theory of coherence capacity. By introducing
capacity currents, transport laws, and focusing theorems, we have shown
that coherence capacity flows, concentrates, and exhausts in a manner
analogous to conserved physical resources.

This dynamical picture completes the structural framework begun in the
preceding work. Gravity, irreversibility, and breakdown of effective laws
emerge not as separate postulates, but as unavoidable consequences of
capacity transport and exhaustion.

\end{document}
